\documentclass[12pt,a4paper]{report}

% ================= PACKAGES =================

\usepackage[a4paper,margin=1in]{geometry}
\usepackage{graphicx}
\usepackage{amsmath,amssymb}
\usepackage{setspace}
\usepackage{xcolor}
\usepackage{hyperref}
\usepackage{fancyhdr}
\usepackage{titlesec}
\usepackage{booktabs}

% ================= FONT =================

\usepackage{newtxtext,newtxmath}

% ================= LINE SPACING =================

\onehalfspacing

% ================= PAGE STYLE =================

\pagestyle{fancy}

\fancyhf{}

\fancyhead[L]{AI for Smart Transportation}
\fancyhead[R]{Master Thesis}

\fancyfoot[C]{\thepage}

% ================= CHAPTER STYLE =================

\titleformat{\chapter}
{\Huge\centering\bfseries\color{black}}
{Chapter \thechapter}{20pt}{}

% ================= DOCUMENT =================

\begin{document}
	
	% =================================================
	% TITLE PAGE
	% =================================================
	
	\begin{titlepage}
		
		\centering
		
		\vspace*{2cm}
		
		{\Huge\bfseries
			Artificial Intelligence for Smart Urban Transportation Systems
			\par}
		
		\vspace{1cm}
		
		{\Large
			A Professional Thesis Template
			\par}
		
		\vspace{2cm}
		
		\rule{6cm}{6cm}
		
		\vspace{2cm}
		
		{\Large
			Brintha Theboral
			\par}
		
		\vspace{1cm}
		
		Department of Computer Science Engineering
		
		Global Institute of Technology
		
		\vfill
		
		Submitted in partial fulfillment of the requirements\\
		for the degree of\\
		Master of Science in Artificial Intelligence
		
		\vfill
		
		{\large May 2026}
		
	\end{titlepage}
	
	% =================================================
	% DECLARATION
	% =================================================
	
	\chapter*{Declaration}
	
	I hereby declare that this thesis entitled
	\textit{Artificial Intelligence for Smart Urban Transportation Systems}
	is an original research work carried out by me.
	
	\newpage
	
	% =================================================
	% ACKNOWLEDGEMENT
	% =================================================
	
	\chapter*{Acknowledgement}
	
	I sincerely thank my supervisor, faculty members, friends, and family members for their valuable support throughout this research work.
	
	\newpage
	
	% =================================================
	% ABSTRACT
	% =================================================
	
	\chapter*{Abstract}
	
	Urban traffic congestion has become a major challenge in modern smart cities. This thesis presents an AI-based framework for intelligent transportation management using deep learning and IoT systems.
	
	The proposed framework integrates real-time sensor analytics, machine learning prediction systems, and adaptive traffic optimization methods to improve transportation efficiency and reduce congestion.
	
	\newpage
	
	% =================================================
	% TABLE OF CONTENTS
	% =================================================
	
	\tableofcontents
	
	\newpage
	
	\listoftables
	
	\newpage
	
	\listoffigures
	
	\newpage
	
	% =================================================
	% CHAPTER 1
	% =================================================
	
	\chapter{Introduction}
	
	\section{Background}
	
	Rapid urbanization has increased transportation challenges worldwide. Smart transportation systems are becoming essential for sustainable urban development.
	
	\section{Problem Statement}
	
	Traditional traffic systems fail to dynamically adapt to changing traffic conditions.
	
	\section{Objectives}
	
	\begin{itemize}
		
		\item Develop AI-based traffic prediction systems
		
		\item Improve congestion management
		
		\item Optimize urban mobility
		
	\end{itemize}
	
	\section{Scope of the Study}
	
	This research focuses on intelligent transportation systems using machine learning and IoT technologies.
	
	% =================================================
	% CHAPTER 2
	% =================================================
	
	\chapter{Literature Review}
	
	\section{Smart Transportation Systems}
	
	Recent advancements in smart city infrastructure have enabled intelligent transportation frameworks.
	
	\section{Machine Learning Approaches}
	
	Several AI techniques have been proposed:
	
	\begin{itemize}
		
		\item Decision Trees
		
		\item Random Forest
		
		\item Deep Neural Networks
		
		\item LSTM Models
		
	\end{itemize}
	
	\section{Research Gap}
	
	Existing systems often lack real-time scalability and adaptive optimization.
	
	% =================================================
	% CHAPTER 3
	% =================================================
	
	\chapter{Methodology}
	
	\section{System Architecture}
	
	The proposed system integrates:
	
	\begin{enumerate}
		
		\item IoT Sensor Networks
		
		\item Data Collection Layer
		
		\item AI Prediction Engine
		
		\item Smart Dashboard
		
	\end{enumerate}
	
	\section{Mathematical Model}
	
	Traffic prediction equation:
	
	\[
	Y_t = f(X_t,W,b)
	\]
	
	Where:
	
	\begin{itemize}
		
		\item $X_t$ = Input traffic data
		
		\item $W$ = Weight parameters
		
		\item $b$ = Bias term
		
	\end{itemize}
	
	Loss function:
	
	\[
	L = \frac{1}{n}\sum_{i=1}^{n}(y_i - \hat{y_i})^2
	\]
	
	\section{Data Collection}
	
	Traffic data was collected from IoT-enabled smart transportation systems.
	
	% =================================================
	% CHAPTER 4
	% =================================================
	
	\chapter{Results and Discussion}
	
	\section{Performance Analysis}
	
	\begin{table}[h]
		
		\centering
		
		\caption{Model Accuracy Comparison}
		
		\begin{tabular}{|c|c|}
			
			\hline
			
			Model & Accuracy \\
			
			\hline
			
			Linear Regression & 78.2\% \\
			
			\hline
			
			Random Forest & 85.6\% \\
			
			\hline
			
			LSTM Network & 93.1\% \\
			
			\hline
			
			CNN-LSTM Hybrid & 95.4\% \\
			
			\hline
			
		\end{tabular}
		
	\end{table}
	
	\section{Discussion}
	
	The hybrid deep learning model achieved superior accuracy compared to traditional machine learning approaches.
	
	\section{Architecture Diagram}
	
	\begin{figure}[h]
		
		\centering
		
		\rule{12cm}{6cm}
		
		\caption{AI-Based Smart Transportation Architecture}
		
	\end{figure}
	
	% =================================================
	% CHAPTER 5
	% =================================================

	\chapter{Conclusion and Future Work}
	
	\section{Conclusion}
	
	This thesis demonstrated the effectiveness of artificial intelligence techniques for smart urban transportation systems.
	
	\section{Future Scope}
	
	Future enhancements include:
	
	\begin{itemize}
		
		\item Edge AI deployment
		
		\item Autonomous vehicle integration
		
		\item Smart traffic signal optimization
		
		\item Real-time analytics
		
	\end{itemize}
	
	% =================================================
	% REFERENCES
	% =================================================
	
	\chapter*{References}
	
	\begin{enumerate}
		
		\item Smith J., Smart Transportation Systems, Springer, 2023.
		
		\item Johnson A., Deep Learning for Traffic Prediction, IEEE Transactions, 2022.
		
		\item Kumar R., AI-Based Urban Analytics, Elsevier, 2021.
		
		\item Lee K., IoT and Smart Mobility Systems, Wiley Publications, 2020.
		
	\end{enumerate}
	
	% =================================================
	% APPENDIX
	% =================================================
	
	\appendix
	
	\chapter{Dataset Information}
	
	The dataset contains:
	
	\begin{itemize}
		
		\item Vehicle count data
		
		\item GPS tracking information
		
		\item Weather conditions
		
		\item Traffic signal timings
		
	\end{itemize}
	
\end{document}